% ============================================================================
%  sections/part3b_deep_layer.tex  --  Part III, deep layer (proofs / subspace /
%  divergence / kernel fixes / numerics).  \input-ready fragment; no preamble.
%  Uses ONLY the macros defined in preamble.tex.
%  Source: notes/07_utility_subspace_proofs_numerics.md (proofs I--III, predictive
%  conditioning, subspace theory, numerical audit).  Ground truth = current code.
% ============================================================================

\section{The deep layer: conditioning, subspace, divergence, and numerics}
\label{sec:deeplayer}

Part~III (\S7) \emph{defined} the two acquisition utilities and their entropies;
this section proves \emph{why} those formulas behave as they do. It supplies the
theorems that sit underneath the moment code, the exact predictive-conditioning
derivation, the subspace-validity argument, the norm-scaling divergence theorems
that explain the utility's blow-up under unconstrained image optimization, the
two kernel-level fixes, and the four-failure-mode numerical audit of the standard
utility.

Two objects recur, and their ``subtracted term'' is \emph{not} the same (the
single most load-bearing clash for a pooled reader; the top-level exposition is
in Part~III):
\begin{itemize}[leftmargin=1.4em]
  \item the \textbf{standard / production} utility
        $\Ustd(x^*)=\Hmarg(x^*)-\E_{\lgr}[\Hnoise(x^*)]$ of \eqref{eq:Ustd},
        which the pool-selection loop maximizes (\S\ref{subsec:numerics} audits
        it); its subtracted term is the expected Poisson \emph{noise} entropy
        $\E_{\lgr}[\Hnoise]$, an average over the point's \emph{own} latent, no
        second image involved;
  \item the \textbf{distribution-aware (DA)} research utility
        $\Uda(x^*)=\Hmarg(x^*)-\E_{x\sim\px}[\Hcond(x^*)]$ of \eqref{eq:Uda},
        whose subtracted term is the entropy \emph{after Gaussian-conditioning}
        the GP on a hypothetical observation of $\lat(x_t)$ at a target image
        $x_t$.
\end{itemize}
The two coincide only in the $\corr^2\!=\!1$ deterministic limit, where $\Hcond$
collapses to a single Poisson entropy (\S\ref{subsec:divergence}). The divergence
theorems of \S\ref{subsec:divergence} and the remedies of
\S\ref{subsec:kernelfixes} concern \emph{gradient ascent on the image $x^*$}
(a super-stimulus / image-synthesis setting), which is distinct from the
production loop's selection from a fixed pool.

% ===========================================================================
\subsection{GP moments and Gaussian conditioning (Proof III)}
\label{subsec:moments}

This is the moments reference grounding \texttt{get\_\allowbreak gp\_\allowbreak marginal\_\allowbreak moments} and
\texttt{get\_\allowbreak gp\_\allowbreak conditional\_\allowbreak moments} (the current helpers; the deprecated
\texttt{lambda\_\allowbreak moments} no longer exists).

\subsubsection{The self-kernel identity}
Recall the order-1 arc-cosine kernel of Part~I, \eqref{eq:acoskernel},
\[
  \Kf(x,x') \;=\; \tfrac1\pi\,\Kmag\,\Jang(\kang),\qquad
  \Kmag=\sqrt{\vmag_{x}\,\vmag_{x'}},\quad
  \cos\kang=\frac{x\transpose\Cmat x'+\sigz^2}{\Kmag},\quad
  \Jang(\kang)=\sin\kang+(\pi-\kang)\cos\kang,
\]
with self-magnitude $\vmag_{x}=x\transpose\Cmat x+\sigz^2$.

\begin{lemma}[Self-kernel / prior variance]\label{lem:selfkernel}
For every $x$, $\Kf(x,x)=\lVert x\rVert_{\Cmat}^2+\sigz^2=\vmag_{x}$, where
$\lVert x\rVert_{\Cmat}^2\equiv x\transpose\Cmat x$.
\end{lemma}
\begin{proof}
Set $x'=x$: $\cos\kang=\vmag_x/\vmag_x=1\Rightarrow\kang=0$, and
$\Jang(0)=\sin0+\pi\cos0=\pi$. Hence
$\Kf(x,x)=\tfrac1\pi\,\vmag_x\,\pi=\vmag_x=x\transpose\Cmat x+\sigz^2$.
\end{proof}

This restates \eqref{eq:selfkernel}: \emph{the prior variance at any point equals
its squared $\Cmat$-norm (plus the bias variance).} It is the single identity
from which the entire norm-scaling divergence of \S\ref{subsec:divergence} grows.

\subsubsection{Variational posterior moments}
With $\Mind$ inducing points $\{\ipt_j\}$ and variational posterior
$q(\latt)=\Gauss(\vm,\vV)$, write the inducing gram $\Ktil=\Kf(\Ipt,\Ipt)$, the
cross-kernel vector $\kvec(x)=[\Kf(x,\ipt_1),\dots,\Kf(x,\ipt_\Mind)]\transpose$,
the projection $\pv(x)=\Ktil^{-1}\kvec(x)$, and the (training-fixed) correction
matrix
\[
  W \;=\; \Ktil^{-1}(\vV-\Ktil)\,\Ktil^{-1}.
\]
The marginal moments are
\begin{equation}\label{eq:margmoments}
\begin{aligned}
  \pmean(x^*) &= \kvec(x^*)\transpose\Ktil^{-1}\vm,\\
  \pvar(x^*)  &= \Kf(x^*,x^*)+\kvec(x^*)\transpose W\,\kvec(x^*)\\
              &= \lVert x^*\rVert_{\Cmat}^2+\sigz^2+\kvec(x^*)\transpose W\,\kvec(x^*).
\end{aligned}
\end{equation}
These ground \texttt{get\_\allowbreak gp\_\allowbreak marginal\_\allowbreak moments}, which
returns \texttt{posterior.\allowbreak mean} and \texttt{posterior.\allowbreak variance}.
The first term of $\pvar$ is the prior variance (Lemma~\ref{lem:selfkernel}); the
second is the posterior correction (negative when $\vV\prec\Ktil$, i.e.\ data
shrink variance). Between two points,
\[
  \Sigma_q(x^*,x_t)=\Kf(x^*,x_t)+\kvec(x^*)\transpose W\,\kvec(x_t),
  \qquad \Sigma_q(x^*,x^*)=\pvar(x^*),\ \ \Sigma_q(x_t,x_t)=\pvar(x_t),
\]
and the squared posterior correlation --- the Pearson correlation of the joint
Gaussian $[\lat(x^*),\lat(x_t)]$ under $q$ --- is
\begin{equation}\label{eq:corrdef}
  \corr^2(x^*,x_t)=\frac{\Sigma_q(x^*,x_t)^2}{\pvar(x^*)\,\pvar(x_t)}\in[0,1].
\end{equation}

\subsubsection{Conditioning on an observed latent}
Observing $\lat_t$ at $x_t$ and applying the standard Gaussian conditioning
identity to $[\lat(x^*),\lat(x_t)]$ gives the conditional moments grounding
\texttt{get\_\allowbreak gp\_\allowbreak conditional\_\allowbreak moments}:
\begin{equation}\label{eq:condmoments}
  \boxed{\;
  \begin{aligned}
  \pmean_{\mathrm{cond}}(x^*)&=\pmean(x^*)+\frac{\Sigma_q(x^*,x_t)}{\pvar(x_t)}\bigl(\lat_t-\pmean(x_t)\bigr),\\[2pt]
  \pvar_{\mathrm{cond}}(x^*)&=\pvar(x^*)-\frac{\Sigma_q(x^*,x_t)^2}{\pvar(x_t)}=\pvar(x^*)\bigl(1-\corr^2\bigr).
  \end{aligned}
  \;}
\end{equation}

\begin{proposition}[$\corr^2$ is the fractional variance reduction]\label{prop:fracvar}
The conditional variance at $x^*$ is the marginal variance scaled by
$(1-\corr^2)$; equivalently, conditioning on $\lat(x_t)$ removes a fraction
$\corr^2$ of the variance at $x^*$.
\end{proposition}
\begin{proof}
Substitute the correlation definition \eqref{eq:corrdef} into the variance line
of \eqref{eq:condmoments}:
\[
  \pvar(x^*)-\frac{\Sigma_q(x^*,x_t)^2}{\pvar(x_t)}
  =\pvar(x^*)\Bigl[1-\frac{\Sigma_q(x^*,x_t)^2}{\pvar(x^*)\,\pvar(x_t)}\Bigr]
  =\pvar(x^*)(1-\corr^2).
\]
\end{proof}

The code (\texttt{get\_\allowbreak gp\_\allowbreak conditional\_\allowbreak moments}) reads the \emph{full joint}
\texttt{covariance\_\allowbreak matrix} of $[x_{\text{sample}};x^*]$ from the model and applies
exactly \eqref{eq:condmoments} (with a floor $\pvar_{\mathrm{cond}}\ge10^{-8}$).
By the theorem of \S\ref{subsec:predictive} this joint route is the \emph{exact}
augmented predictive, not the sparse approximation.

\subsubsection{Log-firing transform and the DA utility in three scalars}
Under $\lgr=\gain\lat+\latz$ and $R\sim\Poi(e^{\lgr})$, the GP moments push to
$\lgr$-space as $\mug=\gain\pmean+\latz$, $\varg=\gain^2\pvar$ (marginal) and
$\mu_{g,\mathrm{cond}}=\gain\,\pmean_{\mathrm{cond}}+\latz$,
$\sigma^{2}_{g,\mathrm{cond}}=\gain^2\pvar_{\mathrm{cond}}=\varg(1-\corr^2)$
(the variance reduction carries through). Hence, writing
$\Hmarg(\mug,\varg)$ for the marginal Poisson--log-normal entropy \eqref{eq:Hmarg},
\[
  \Uda(x^*)=\Hmarg(\mug,\varg)-\Hmarg\!\left(\mu_{g,\mathrm{cond}},\,\sigma^{2}_{g,\mathrm{cond}}\right).
\]
In deterministic mode ($\lat_t=\pmean(x_t)$, \texttt{sample\_\allowbreak lambda=False}) the
conditional mean is unchanged, $\mu_{g,\mathrm{cond}}=\mug$, and \eqref{eq:Uda}
reduces to the deterministic distribution-aware utility \eqref{eq:da-deterministic}
of \S\ref{subsubsec:det}, a function of exactly three scalars:
\begin{equation}\label{eq:daboxed}
  \boxed{\;\Uda(x^*)=\Hmarg(\mug,\varg)-\Hmarg\!\left(\mug,\ \varg(1-\corr^2)\right)\;}
\end{equation}
the operating point $\mug=\gain\pmean(x^*)+\latz$, the marginal variance
$\varg=\gain^2\pvar(x^*)$, and the fractional reduction $\corr^2$. Conditioning
slides $(\mug,\varg)\!\to\!(\mug,\varg(1-\corr^2))$: same mean, reduced variance;
the utility is the entropy gap. This is the object whose divergence
\S\ref{subsec:divergence} analyses.

% ===========================================================================
\subsection{The predictive distribution conditioned on a new observation}
\label{subsec:predictive}

The DA utility integrates the entropy of $p(\lat(x^*)\mid\lat(x),\mathcal{D})$:
the predictive over the latent at a query $x^*$ after a \emph{hypothetical} new
latent observation at $x$. We derive it two ways --- sparse-approximated and
exact/augmented --- and prove the two agree in mean.

Define the \emph{prior conditional covariance} (the residual not captured by the
inducing points)
\[
  \ccond(x,x') \;\equiv\; \Kf(x,x')-\pv(x)\transpose\Ktil\,\pv(x')
                        \;=\; \Kf(x,x')-\kvec(x)\transpose\Ktil^{-1}\kvec(x'),
\]
whose diagonal is the residual / aleatoric variance
$\varr(x)\equiv\ccond(x,x)=\Kf(x,x)-\kvec(x)\transpose\Ktil^{-1}\kvec(x)\ge0$.

\subsubsection{The rank-1 posterior update}
\begin{lemma}[Rank-1 update]\label{lem:rank1}
Let the new point $x$ have projection $\pv=\pv(x)$ and residual variance
$\varr=\varr(x)$, so the conditional likelihood is
$p(\lat(x)\mid\latt)=\Gauss(\pv\transpose\latt,\varr)$. Conditioning the
variational posterior $\Gauss(\vm,\vV)$ on the observed $\lat(x)$ yields
$\Gauss(\vm',\vV')$ with
\begin{equation}\label{eq:rank1}
  \vm' \;=\; \vm+\frac{\vV\pv}{\varr+\pv\transpose\vV\pv}\bigl(\lat(x)-\pv\transpose\vm\bigr),
  \qquad
  \vV' \;=\; \vV-\frac{\vV\pv\,\pv\transpose\vV}{\varr+\pv\transpose\vV\pv}.
\end{equation}
\end{lemma}
\begin{proof}
Under the variational approximation the joint of $(\lat(x),\latt)$ is Gaussian:
\[
  \begin{bmatrix}\lat(x)\\ \latt\end{bmatrix}
  \sim\Gauss\!\left(
  \begin{bmatrix}\pv\transpose\vm\\ \vm\end{bmatrix},\
  \begin{bmatrix}\varr+\pv\transpose\vV\pv & \pv\transpose\vV\\[2pt]
                 \vV\pv & \vV\end{bmatrix}\right).
\]
The Gaussian conditioning identity on $\latt$ given $\lat(x)$, with scalar
lower-right-to-upper-left Schur term $\Sigma_{11}=\varr+\pv\transpose\vV\pv$ and
cross-block $\vV\pv$, gives \eqref{eq:rank1} directly.
\end{proof}

\subsubsection{Sparse vs.\ augmented predictive}
Let $\pv^*=\Ktil^{-1}\kvec(x^*)$ and $\varr(x^*)=\Kf(x^*,x^*)-\kvec(x^*)\transpose\Ktil^{-1}\kvec(x^*)$.

\emph{Sparse route.} Assuming conditional independence $p(\lat(x^*)\mid\lat(x),\latt)\approx p(\lat(x^*)\mid\latt)$ and integrating the conditional prior against the updated posterior of Lemma~\ref{lem:rank1},
\[
  \pmean_{\mathrm{pred}}={\pv^*}\transpose\vm',\qquad
  \pvar_{\mathrm{pred}}=\varr(x^*)+{\pv^*}\transpose\vV'\pv^*
                       =\Kf(x^*,x^*)-{\pv^*}\transpose(\Ktil-\vV')\pv^*.
\]

\emph{Augmented (exact) route.} Keep the direct $\lat(x)$--$\lat(x^*)$
correlation. With augmented latent $\latt_{\mathrm{aug}}=[\lat(x);\latt]$ and
augmented gram / cross-kernel
\[
  \widetilde{K}_{x}=\begin{bmatrix}\Kf(x,x) & \kvec(x)\transpose\\ \kvec(x) & \Ktil\end{bmatrix},
  \qquad
  \kvec_{\mathrm{aug}}(x^*)=\begin{bmatrix}\Kf(x,x^*)\\ \kvec(x^*)\end{bmatrix},
\]
the non-approximated conditional prior is
$p(\lat(x^*)\mid\latt_{\mathrm{aug}})=\Gauss(\pv_{\mathrm{aug}}\transpose\latt_{\mathrm{aug}},S^*)$,
$\pv_{\mathrm{aug}}=\widetilde{K}_x^{-1}\kvec_{\mathrm{aug}}(x^*)$,
$S^*=\Kf(x^*,x^*)-\pv_{\mathrm{aug}}\transpose\widetilde{K}_x\,\pv_{\mathrm{aug}}$.
Since $\lat(x)$ is \emph{observed} (variance $0$) while $\latt$ follows the
updated posterior $\Gauss(\vm',\vV')$, the augmented posterior covariance is
$V''=\bigl[\begin{smallmatrix}0&0\transpose\\0&\vV'\end{smallmatrix}\bigr]$, and
\[
  \pmean_{\mathrm{pred}}=\pv_{\mathrm{aug}}\transpose[\lat(x);\vm'],\qquad
  \pvar_{\mathrm{pred}}=S^*+\pv_{\mathrm{aug}}\transpose V''\pv_{\mathrm{aug}}
                       =\Kf(x^*,x^*)-\pv_{\mathrm{aug}}\transpose(\widetilde{K}_x-V'')\,\pv_{\mathrm{aug}}.
\]

\begin{remark}[Augmented route fixes the sparse self-variance artifact]
When $x=x^*$, $\pv_{\mathrm{aug}}$ aligns with the first column of
$\widetilde{K}_x$ and, because the $(1,1)$ block of $V''$ is $0$, the predictive
variance is forced to $0$. A self-consistent predictive \emph{must} have zero
variance at an already-observed point; the sparse route leaves a spurious nonzero
self-variance there. This is the same pathology as the hand-assembled
cross-covariance that can violate Cauchy--Schwarz (Part~II \S4.7) --- the shipped
code sidesteps both by reading the model's full joint covariance.
\end{remark}

\subsubsection{Equivalence of predictive means}
\begin{theorem}[Equivalence of predictive means]\label{thm:predmean-equiv}
The augmented-integration predictive mean equals the mean obtained by standard
Gaussian conditioning on the joint predictive $[\lat(x),\lat(x^*)]$ --- exactly,
with no cross-covariance approximation.
\end{theorem}
\begin{proof}
\textbf{Method 1 (joint conditioning).} The joint moments given $\mathcal{D}$ are
$\pmean_x=\pv\transpose\vm$, $\pmean_{x^*}={\pv^*}\transpose\vm$,
\[
  \Sigma_{11}=\varr+\pv\transpose\vV\pv,\qquad
  \Sigma_{12}=\underbrace{\bigl(\Kf(x,x^*)-\pv\transpose\Ktil\,\pv^*\bigr)}_{=\ \ccond}+{\pv^*}\transpose\vV\pv,
\]
where the residual cross term is exactly $\ccond=\ccond(x,x^*)$. The Gaussian
conditioning identity gives the target
\begin{equation}\label{eq:method1}
  \E[\lat(x^*)\mid\lat(x),\mathcal{D}]
  =\pmean_{x^*}+\frac{\Sigma_{12}}{\Sigma_{11}}\bigl(\lat(x)-\pmean_x\bigr)
  ={\pv^*}\transpose\vm+\frac{\Sigma_{12}}{\Sigma_{11}}\bigl(\lat(x)-\pv\transpose\vm\bigr).
\end{equation}

\textbf{Method 2 (augmented integration).} Block-invert $\widetilde{K}_x$. Writing
the residual-correlation coefficient $\vartheta\equiv\ccond/\varr$ (the source's
predictive $\alpha$, renamed here to avoid the RF locality mask $\loc_i$),
\[
  \pv_{\mathrm{aug}}=\begin{bmatrix}\vartheta\\ \pv^*-\vartheta\,\pv\end{bmatrix},
  \qquad \vartheta=\frac{\Kf(x,x^*)-\pv\transpose\Ktil\,\pv^*}{\varr}=\frac{\ccond}{\varr}.
\]
Then, with $\vm'$ from \eqref{eq:rank1},
\[
  \pmean_{\mathrm{aug}}=\pv_{\mathrm{aug}}\transpose[\lat(x);\vm']
  =\vartheta\,\lat(x)+({\pv^*}\transpose-\vartheta\,\pv\transpose)\vm'
  =\vartheta\,\lat(x)+({\pv^*}\transpose-\vartheta\,\pv\transpose)\!\left[\vm+\frac{\vV\pv}{\Sigma_{11}}\bigl(\lat(x)-\pv\transpose\vm\bigr)\right].
\]
Group by the prior mean ${\pv^*}\transpose\vm$ and the innovation
$(\lat(x)-\pv\transpose\vm)$, adding and subtracting $\vartheta\,\pv\transpose\vm$:
\[
  \pmean_{\mathrm{aug}}={\pv^*}\transpose\vm
   +\underbrace{\left[\vartheta+\frac{{\pv^*}\transpose\vV\pv-\vartheta\,\pv\transpose\vV\pv}{\Sigma_{11}}\right]}_{\displaystyle \mathcal{C}}
   \bigl(\lat(x)-\pv\transpose\vm\bigr).
\]
Put $\mathcal{C}$ over the common denominator $\Sigma_{11}=\varr+\pv\transpose\vV\pv$
and use $\vartheta\,\varr=\ccond$:
\[
  \mathcal{C}=\frac{\vartheta\,\varr+\vartheta\,\pv\transpose\vV\pv+{\pv^*}\transpose\vV\pv-\vartheta\,\pv\transpose\vV\pv}{\Sigma_{11}}
   =\frac{\vartheta\,\varr+{\pv^*}\transpose\vV\pv}{\Sigma_{11}}
   =\frac{\ccond+{\pv^*}\transpose\vV\pv}{\Sigma_{11}}
   =\frac{\Sigma_{12}}{\Sigma_{11}}.
\]
Hence $\pmean_{\mathrm{aug}}={\pv^*}\transpose\vm+(\Sigma_{12}/\Sigma_{11})(\lat(x)-\pv\transpose\vm)$,
identical to \eqref{eq:method1}.
\end{proof}

\begin{remark}
Theorem~\ref{thm:predmean-equiv} certifies that the production conditional-moment
code (\S\ref{subsec:moments}, \texttt{get\_\allowbreak gp\_\allowbreak conditional\_\allowbreak moments}), which
conditions on the full joint covariance, computes the \emph{exact} augmented
predictive mean --- it does not incur the sparse cross-covariance error.
Setting the residual cross term $\ccond=0$ would enforce the sparse factorization
\[
  p(\lat^*,\lat\mid\latt)\approx p(\lat^*\mid\latt)\,p(\lat\mid\latt);
\]
the inducing (epistemic) cross-covariance ${\pv^*}\transpose\vV\pv$ can still be
nonzero, because changing belief about an inducing point moves both $\lat(x^*)$
and $\lat(x)$.
\end{remark}

% ===========================================================================
\subsection{Subspace theory}
\label{subsec:subspace}

\emph{Why} the utility may be optimized in a low-dimensional subspace of image
space (11664 pixels full; $\sim$900--2356 in the RF-masked region) instead of
full pixel space. Setting: optimize $x^*$ to maximize $\Uda(x^*)$ by gradient
ascent, constrained to $x_{\mathrm{rf}}=\bar{x}+\text{(basis)}\cdot\subc$ over the
subspace coordinate $\subc$; the offset $\bar{x}$ is the RF-pixel dataset mean, so
$\subc=0$ maps to the mean image.

\subsubsection{The kernel touches $x$ only through the $\Cmat$-eigenspace}
Let $\Cmat=\sum_i\gamma_i\,e_i e_i\transpose$ be the eigendecomposition of the
structured PSD matrix of Part~I (\S3.3), with $\gamma_1\ge\dots\ge\gamma_d\ge0$
and orthonormal $\Cmat$-eigenvectors $e_i$.

\begin{proposition}[$\Cmat$-eigenspace validity]\label{prop:subspace}
Every kernel evaluation depends on $x$ only through the quadratic forms
$x\transpose\Cmat x=\sum_i\gamma_i(e_i\transpose x)^2$ and $x\transpose\Cmat x'$;
and $\nabla_x\Kf$ lies in the column space of $\Cmat$, with its component along
$e_i$ scaled by $\gamma_i$. Consequently, gradient ascent already concentrates
movement along the top $\Cmat$-eigenvectors, and truncating to the leading
$\Cmat$-eigen-subspace is (at full rank) an \emph{exact reparameterization} and
(at reduced rank) drops only near-zero-gradient directions.
\end{proposition}
\begin{proof}
The kernel enters $x$ solely via $x\transpose\Cmat x+\sigz^2=\vmag_x$ and
$x\transpose\Cmat x'+\sigz^2$ (through $\Kmag$ and $\cos\kang$); in the
$\Cmat$-eigenbasis $x\transpose\Cmat x=\sum_i\gamma_i(e_i\transpose x)^2$, so a
direction $e_i$ with $\gamma_i\approx0$ contributes $\approx0$ to \emph{every}
kernel value. The input gradient of the arc-cosine kernel (Part~I \S3.5,
\eqref{eq:gradxK}) is
\[
  \nabla_x\Kf(x,x')=\tfrac1\pi\Bigl[(\pi-\kang)\,\Cmat x'+\sin\kang\,\sqrt{\vmag_{x'}/\vmag_x}\;\Cmat x\Bigr],
\]
every term of which is $\Cmat\cdot(\text{vector})$; projecting onto $e_i$ multiplies
by $\gamma_i$. Thus movement along small-$\gamma_i$ directions is intrinsically
suppressed. Removing those directions (a) does not change the optimum (they carried
near-zero gradient) and (b) removes flat directions, giving a better-conditioned,
lower-effective-dimension problem.
\end{proof}

So the $\Cmat$-eigenspace projection is a \emph{convergence aid / exact
reparameterization}, not a naturalness constraint (analogous to the eigenvalue
threshold on $\Ktil$ in the direct-VGP path of Part~II \S5.1).

\subsubsection{Three bases and the Szeg\H{o}/Fourier equivalence}
All three write $x_{\mathrm{rf}}=\bar{x}+\text{(basis)}\cdot\subc$ and share one
optimizer/diagnostics; they differ only in the basis.
\begin{itemize}[leftmargin=1.4em]
  \item \textbf{PCA.} Columns $\Phi_K=[\phi_1,\dots,\phi_K]$ = top-$K$ eigenvectors
        of the data covariance $\Sigma_{\mathrm{data}}=\tfrac1N X_c\transpose X_c$
        (eigenvalues $\nu_1\ge\nu_2\ge\dots$). The offset $\bar{x}$ is
        \emph{mathematically intrinsic} (centering before SVD); truncation is a
        \emph{genuine data-driven constraint} (excludes low-data-variance,
        ``unnatural'' directions). At variance threshold $0.80$: $K=197$ of
        $2356$ RF pixels. Gradient projects as $\nabla_\subc\Uda=\Phi_K\transpose\nabla_x\Uda$.
        PCA captures \emph{second-order statistics only} (mean $+$ covariance);
        it is blind to phase structure (edges/contours), sparsity, and
        non-Gaussian marginals.
  \item \textbf{$\Cmat$-eigenspace.} Columns $E_K=[e_1,\dots,e_K]$ = top-$K$
        $\Cmat$-eigenvectors. The offset is \emph{not} intrinsic (the
        eigendecomposition passes through the origin) --- see the caveat below.
        Truncation is the reparameterization of Proposition~\ref{prop:subspace}.
        The spectrum is extremely skewed (top eigenvalue $\sim44\%$ of the total):
        at relative threshold $10^{-3}$, $K=28$ of $2356$ captures $99.7\%$ of the
        eigenvalue mass. (Float32 noise floor: eigenvalues below $\sim10^{-7}\gamma_1$
        are numerical noise; thresholds below $\sim10^{-6}$ pull in garbage
        eigenvectors.)
  \item \textbf{Combined.} Project $\Cmat$ into the PCA subspace and re-diagonalize:
        $\Cmat_{\mathrm{pca}}=\Phi_K\transpose\Cmat\,\Phi_K=Q\,\Gamma_Q Q\transpose$,
        keep the columns $Q$ above the eigen-threshold, basis $=\Phi_K Q$.
        Directions that are \emph{both} natural (PCA) and kernel-visible
        ($\Cmat$-eigen). At $(0.80,10^{-3})$: $2356\to K_{\mathrm{pca}}=197\to
        K_{\mathrm{comb}}=28$, achieving near-identical utility to PCA with
        $\sim7\times$ fewer dimensions.
\end{itemize}

\begin{proposition}[PCA $\approx$ Fourier under stationarity (Szeg\H{o})]
\label{prop:szego}
If the image process is second-order stationary on a regular grid, its covariance
is (block-)Toeplitz, $\Sigma[i,j]=f(|i-j|)$. By Szeg\H{o}'s theorem, as $d\to\infty$
the eigenvectors converge to Fourier modes and the eigenvalues to the power
spectral density $S(\omega)=\sum_k f(k)e^{-i\omega k}$. Natural images have
$S(\omega)\sim1/|\omega|^2$, so $\Sigma_{\mathrm{data}}$'s eigenvectors approach
the Fourier basis and PCA-truncation approximates keeping the $K$
lowest-frequency Fourier modes.
\end{proposition}
The equivalence is weakened by the irregular masked region, the non-stationarity
the $\loc$ mask injects (RF-center pixels have higher variance, breaking Toeplitz
structure), finite $N\approx10^4$ in $d\approx900$, and residual
non-stationarity of natural images; at $\sim30\times30$ patches it is reasonable.

\subsubsection{Subspace $\ne$ naturalness constraint}
All three are \emph{subspace} constraints, not \emph{distributional} ones. None
guarantees $x^*$ is a plausible sample of $\px$: an image with the right power
spectrum but wrong phase (correlated noise) satisfies all three. A true
distributional constraint needs higher-order structure (phases, edges, sparsity)
--- i.e.\ a generative model, which is the motivation for the guided-diffusion
synthesis of Part~IV. The subspace approach is a useful \emph{diagnostic} (does
restricting to the right power spectrum change utility? does the GP utility
operate primarily at second order?), not a full fix.

\caveat{The $\Cmat$-eigenspace offset is not mathematically motivated. Unlike PCA,
where the mean offset is intrinsic to centering, the $\Cmat$-eigendecomposition
passes through the origin, so setting $\bar{x}=\text{mean}$ has no principled
basis --- it is forced by \emph{limiter precision}: the non-capturable
eigenvectors collapse to $0$ instead of to the mean, and the code compensates by
setting offset $=$ mean. This changes the utility landscape (the kernel sees
$\Cmat(\bar{x}+E_K\subc)$, not $\Cmat(E_K\subc)$) and should be investigated. It
is an acknowledged modeling wrinkle, not a resolved result.}

% ===========================================================================
\subsection{Divergence theorems (Proof I): norm scaling and the utility blow-up}
\label{subsec:divergence}

The arc-cosine kernel is positively homogeneous; scaling an image along a fixed
direction inflates the posterior variance as the square of the scale while
preserving the conditioning benefit, so unconstrained gradient ascent on the DA
utility diverges. The three source proof documents each state this phenomenon
independently; it is proved once here and reused by \S\ref{subsec:kernelfixes}.
Throughout, $x^*=\scal\,x_t$ with scale $\scal>0$, and $q\equiv x_t\transpose\Cmat x_t$.

\begin{theorem}[Cross-kernel proportionality, $\sigz^2=0$]\label{thm:divergence}
Let $\sigz^2=0$ and $x^*=\scal\,x_t$, $\scal>0$. Then $\Kf(x^*,z)=\scal\,\Kf(x_t,z)$
for every $z$, and hence $\kvec(x^*)=\scal\,\kvec(x_t)$.
\end{theorem}
\begin{proof}
With $\sigz^2=0$, $\vmag_{x^*}=\scal^2 q$ and $\vmag_{x_t}=q$. For any $z$ with
$\vmag_z=z\transpose\Cmat z$,
\[
  \cos\kang(x^*,z)=\frac{(\scal x_t)\transpose\Cmat z}{\sqrt{\scal^2 q\cdot\vmag_z}}
   =\frac{\scal\,x_t\transpose\Cmat z}{\scal\sqrt{q\,\vmag_z}}
   =\frac{x_t\transpose\Cmat z}{\sqrt{q\,\vmag_z}}=\cos\kang(x_t,z),
\]
so $\kang(x^*,z)=\kang(x_t,z)$ and $\Jang(\kang(x^*,z))=\Jang(\kang(x_t,z))$. Then
$\Kf(x^*,z)=\tfrac1\pi\sqrt{\scal^2 q\,\vmag_z}\,\Jang(\kang)=\scal\cdot\tfrac1\pi\sqrt{q\,\vmag_z}\,\Jang(\kang)=\scal\,\Kf(x_t,z)$.
Applying this to each inducing point $z=\ipt_j$ gives $\kvec(x^*)=\scal\,\kvec(x_t)$.
\end{proof}

\begin{corollary}[Perfect posterior correlation]\label{cor:rho1}
Under Theorem~\ref{thm:divergence}, $\corr^2(x^*,x_t)=1$.
\end{corollary}
\begin{proof}
From $\kvec(x^*)=\scal\,\kvec(x_t)$ and (with $x^*\parallel x_t$, $\kang=0$,
$\Jang(0)=\pi$) $\Kf(x^*,x_t)=\scal\,\Kf(x_t,x_t)=\scal\,\vmag_{x_t}$, substitute
into $\Sigma_q$ with $W=\Ktil^{-1}(\vV-\Ktil)\Ktil^{-1}$:
\[
  \Sigma_q(x^*,x_t)=\scal\,\vmag_{x_t}+\scal\,\kvec(x_t)\transpose W\kvec(x_t)=\scal\,\Sigma_q(x_t,x_t),
  \quad
  \Sigma_q(x^*,x^*)=\scal^2\,\Sigma_q(x_t,x_t),
\]
so $\corr^2=[\scal\,\Sigma_q(x_t,x_t)]^2/[\scal^2\Sigma_q(x_t,x_t)\cdot\Sigma_q(x_t,x_t)]=1$.
\end{proof}

\begin{corollary}[Moment scaling]\label{cor:momentscale}
Under Theorem~\ref{thm:divergence}, $\pmean(\scal x_t)=\scal\,\pmean(x_t)$ and
$\pvar(\scal x_t)=\scal^2\,\pvar(x_t)$.
\end{corollary}
\begin{proof}
Substitute $\kvec(x^*)=\scal\,\kvec(x_t)$ into $\pmean=\kvec\transpose\Ktil^{-1}\vm$
(linear in $\kvec$) and $\pvar=\Kf(x^*,x^*)+\kvec\transpose W\kvec$ (quadratic in
$\kvec$, with $\Kf(x^*,x^*)=\scal^2\vmag_{x_t}$).
\end{proof}

By Corollary~\ref{cor:rho1}, $\corr^2=1\Rightarrow\pvar_{\mathrm{cond}}(x^*)=0$:
conditioning removes \emph{all} the variance. Yet by
Corollary~\ref{cor:momentscale} the marginal moments still grow
($\mug\propto\scal$, $\varg\propto\scal^2\Rightarrow\Hmarg\to\infty$), so $\Uda$
still diverges. In the exact parallel case ($\kang=0,\sigz^2=0$),
$\pmean_{\mathrm{cond}}=\pmean$, $\pvar_{\mathrm{cond}}=0$, and $\Hcond$ collapses
to a single Poisson entropy $H_{\Poi}(e^{\mug})$, $\mug=\gain\,\pmean(x^*)+\latz$;
then \eqref{eq:daboxed} becomes
$\Uda(\scal x_t)=\Hmarg(\mug,\varg)-H_{\Poi}(e^{\mug})$ with
$\mug=\gain\scal\,\pmean(x_t)+\latz$, $\varg=\gain^2\scal^2\,\pvar(x_t)$.

\subsubsection{Growth rate of the diverging utility}
\begin{proposition}[Gaussian observation model: $\log\scal$ growth]\label{prop:gauss-logc}
If observations were Gaussian, $r(x)\sim\Gauss(\lat(x),\tau^2)$ with fixed $\tau^2$,
then for $\corr^2=1$, $\pvar=\scal^2\pvar(x_t)$, the DA utility grows as $\log\scal$.
\end{proposition}
\begin{proof}
$\Hmarg=\tfrac12\log(2\pi e(\pvar+\tau^2))$ and, since $\pvar_{\mathrm{cond}}=0$,
$\Hcond=\tfrac12\log(2\pi e\,\tau^2)$, so
$\Uda=\tfrac12\log(1+\pvar/\tau^2)$. With $\pvar=\scal^2\pvar(x_t)$ and large
$\scal$, $\Uda\approx\tfrac12\log(\scal^2\pvar(x_t)/\tau^2)=\log\scal+\text{const}$.
\end{proof}
Unbounded but slow: $d\Uda/d\scal\sim1/\scal$, so ascent slows but never stops.
Norm scaling is \emph{not} unique to Poisson --- the distinction is the \emph{rate}.

\begin{proposition}[Poisson-GP: $\scal^2$ heuristic bound]\label{prop:poisson-c2}
For the Poisson-GP case, under the heuristic $\Hmarg\ge H_{\Poi}(\E[f])$ with
$\E[f]=e^{\mug+\varg/2}$ and the large-rate Poisson entropy
$H_{\Poi}(\lambda)\approx\tfrac12\log(2\pi e\lambda)$,
\[
  \Hmarg\ge\tfrac12\log(2\pi e)+\tfrac12(\mug+\varg/2),
  \qquad\text{dominant term }\tfrac14\varg=\tfrac14\gain^2\scal^2\,\pvar(x_t).
\]
Since $\Hcond=H_{\Poi}(e^{\mug})$ depends only on $\mug$ (grows as $\scal$, not
$\scal^2$),
\[
  \Uda=\Hmarg-\Hcond\ \gtrsim\ \tfrac14\varg\ \propto\ \scal^2
\]
(faster than the Gaussian $\log\scal$).
\end{proposition}
\caveat{The bound $\Hmarg\ge H_{\Poi}(\E[f])$ is \emph{plausible but not proved}.
A Poisson mixture over log-normal rates is overdispersed (``more spread''), but
the Poisson is \emph{not} the max-entropy law on $\mathbb{N}_0$ under a mean
constraint alone (the geometric is); it is max-entropy only under $\mathrm{Var}=\text{mean}$,
so the naive argument does not directly apply. Numerically (\texttt{verify\_\allowbreak conclusions.py},
test C4), fitting $\Uda\sim\scal^\alpha$ over the Laplace-valid range
$\scal\in\{2,5,10\}$ (all $\sum p(r)>0.95$) gives $\alpha\approx1.9$, consistent
with $\scal^2$. The Laplace approximation truncates at $\rmax$ and breaks once
$\varg$ puts mass above $\rmax$; the safety score
$z_{\mathrm{safe}}=(\log\rmax-\mug)/\sqrt{\varg}$ falls below $2$ (unreliable) at
$\scal\approx10$ in-model, limiting the verifiable range.}

\subsubsection{The $\sigz^2>0$ damping and the alignment peak at $\scal=1$}
With $\sigz^2>0$ the results hold only approximately:
\[
  \vmag_{\scal x_t}=\scal^2 q+\sigz^2\neq\scal^2\vmag_{x_t},\qquad
  \cos\kang(\scal x_t,z)=\frac{\scal\,x_t\transpose\Cmat z+\sigz^2}{\sqrt{(\scal^2 q+\sigz^2)\,\vmag_z}}\neq\cos\kang(x_t,z).
\]
The fractional error $|\Kf(\scal x,z)-\scal\Kf(x,z)|/|\scal\Kf(x,z)|$ converges to a
\emph{constant} $O(\sigz^2/\vmag_{x_t})$ as $\scal\to\infty$ (not to zero: $\sigz^2$
shifts the angle by an amount $\propto\sigz^2/\vmag_{x_t}$, independent of $\scal$).
In a trained model ($\vmag_{x_t}=q+\sigz^2\approx669$, $\sigz^2\approx0.98$, so
$\sigz^2/\vmag_{x_t}\approx0.0015$): the measured fractional error in
$\Kf(\scal x_t,\ipt_j)\to\approx0.6\%$ for large $\scal$, $\corr^2\to\approx0.987$,
and the residual variance ratio $\approx1.3\%$ (conditioning stays very effective,
not perfect). Crucially, the cosine similarity
\[
  \cos\kang(\scal)=\frac{\scal q+\sigz^2}{\sqrt{(\scal^2 q+\sigz^2)(q+\sigz^2)}}
\]
is \emph{strictly maximized at $\scal=1$} when $\sigz^2>0$ (and $\equiv1$ for all
$\scal$ when $\sigz^2=0$). So $\sigz^2$ breaks scale invariance: the kernel
\emph{wants} to match the target in pattern \emph{and} intensity --- but the
magnitude still scales as $\Kf(x^*,x^*)=\scal^2 q+\sigz^2=O(\scal^2)$, so the
bounded correlation reward cannot compete with the unbounded variance reward, and
ascent still runs $\scal\to\infty$.

\begin{remark}[Root cause and remedy]
The mechanism is a decoupling of norm from direction: $\Kf(\scal x,y)=\scal\Kf(x,y)$
at $\sigz^2=0$ lets ascent raise variance (hence entropy and utility) via the norm
degree of freedom without sacrificing directional alignment. Any practical
gradient-based image optimization with this kernel needs a norm constraint or a
normalized kernel (\S\ref{subsec:kernelfixes}). The blow-up is strongly
angle-dependent (test C8): effective at small angles ($\kang<0.2$, strong
conditioning), negligible at large angles ($\kang>0.6$).
\end{remark}

Table~\ref{tab:divevidence} shows the trained-model evidence: the DA-optimized
image finds a \emph{genuine} loophole ($\Uda=0.800$ at a small angle
$\kang=0.114$, variance ratio $0.060\Rightarrow\corr^2\approx0.94$), while the
standard-optimized image diverges to $\Ustd=\infty$ by numerical
overflow at $\mug=1816$ --- two different phenomena
(\S\ref{subsec:numerics} separates them).

\begin{table}[t]\centering\small
\caption{Trained model (seed $42$, cell $8$, $\Mind=50$, $\Ntr=50$; $\gain=0.0265$,
$\latz=-1.686$). The small gain $\gain$ compresses the raw GP mean into the useful
band, so the DA loophole ($\Uda=0.800$) is real, not a numerical artifact; the
standard-optimized point overflows float32.}
\label{tab:divevidence}
\begin{tabular}{lrrrrrrr}
\toprule
Image & $\pmean$ & $\pvar$ & $\mug$ & $\varg$ & $\kang$(rad) & $\Hmarg$ & $\Uda$\\
\midrule
$x_t$ (natural target) & $7.86$ & $70.9$ & $-1.48$ & $0.050$ & $0$ & $0.593$ & $0.010$\\
$x^*_{\mathrm{DA}}$ & $101$ & $3240$ & $0.99$ & $2.28$ & $0.114$ & $2.836$ & $0.800$\\
$x^*_{\mathrm{Std}}$ & $68497$ & $2.85\!\times\!10^{8}$ & $1816$ & $2.0\!\times\!10^{5}$ & $0.866$ & $\sim0$ & $\infty$\\
\bottomrule
\end{tabular}
\end{table}

% ===========================================================================
\subsection{Kernel solutions (Proof II)}
\label{subsec:kernelfixes}

Two kernel-level remedies bound the variance term that
\S\ref{subsec:divergence} showed to be the culprit. Decompose the utility (Poisson
link, $\rate=e^{\lat}$) schematically as
\[
  \Ustd(x^*)\ \approx\ \underbrace{\tfrac12\pvar(x^*)}_{\text{variance term, }O(\scal^2)}\ +\ \underbrace{I(\corr^2)}_{\text{correlation term, }\le\text{const}} ,
\]
where the correlation term is bounded (maximized at $\scal=1$ once $\sigz^2>0$) but
the variance term is unbounded; to fix the divergence, bound the variance term.

\subsubsection{Normalized arc-cosine kernel (project onto the unit sphere)}
\begin{proposition}[Scale-invariant normalized kernel]\label{prop:normkernel}
Define $\bar{\Kf}(x,x')=\Kf(x,x')/\sqrt{\Kf(x,x)\,\Kf(x',x')}$. For the order-1
arc-cosine kernel, $\bar{\Kf}(x,x')=\Jang(\kang)/\pi$, and for any $\scal>0$,
$\bar{\Kf}(\scal x,\scal x)=1=\bar{\Kf}(x,x)$ --- the prior variance is strictly
constant, independent of $\lVert x\rVert$.
\end{proposition}
\begin{proof}
By Lemma~\ref{lem:selfkernel} $\Kf(x,x)=\vmag_x$, so
$\bar{\Kf}(x,x')=[\tfrac1\pi\sqrt{\vmag_x\vmag_{x'}}\,\Jang(\kang)]/\sqrt{\vmag_x\vmag_{x'}}=\Jang(\kang)/\pi$,
which depends only on $\kang$. In particular $\bar{\Kf}(x,x)=\Jang(0)/\pi=1$; and
$\bar{\Kf}(\scal x,\scal x)=\Kf(\scal x,\scal x)/\sqrt{\Kf(\scal x,\scal x)^2}=1$.
\end{proof}
With $\bar{\Kf}(x,x)=1$ and cross terms $\bar{\Kf}(x,\ipt_j)=\Jang(\kang_j)/\pi\in[0,1]$
bounded, the whole posterior-correction term is bounded, so $\bar{\pvar}(x)$ is
bounded regardless of $\lVert x\rVert_{\Cmat}$ and $\Uda$ cannot diverge by
scaling: the only way up is to improve $\corr^2$ (small angle to $x_t$; with
$\sigz^2>0$ retained inside $\kang$, minimized only when $x\to x_t$ in both
direction and magnitude).
\caveat{The normalized kernel discards all norm information, but image norm carries
genuine neural-encoding signal (neurons respond differently to high/low contrast).
Empirically, training with $\bar{\Kf}$ drops test correlation by $\sim25\%$ --- the
norm encodes real structure the GP needs. A fix for the image-synthesis divergence
is not a free improvement to the production model.}

\subsubsection{Biologically plausible saturation kernel (arc-sin)}
The arc-cosine kernel is the covariance of infinite ReLU networks, which do
\emph{not} saturate; real neurons have refractory periods and maximum firing
rates. The arc-sin (probit-activation) kernel models saturation:
\[
  \Kf_{\mathrm{sat}}(x,x')=\satf\cdot\tfrac2\pi\,\arcsin\!\left(\frac{x\transpose\Cmat x'+\sigz^2}{\sqrt{(1+x\transpose\Cmat x+\sigz^2)(1+{x'}\transpose\Cmat x'+\sigz^2)}}\right).
\]
\begin{proposition}[Bounded self-variance]\label{prop:satkernel}
$\displaystyle\lim_{\lVert x\rVert\to\infty}\Kf_{\mathrm{sat}}(x,x)=\satf\cdot\tfrac2\pi\arcsin(1)=\satf$.
\end{proposition}
\begin{proof}
As $\lVert x\rVert\to\infty$ the ratio inside $\arcsin$ tends to
$(x\transpose\Cmat x)/(1+x\transpose\Cmat x)\to1$, and $\arcsin(1)=\pi/2$, so
$\Kf_{\mathrm{sat}}(x,x)\to\satf\cdot(2/\pi)(\pi/2)=\satf$.
\end{proof}
Unlike $\Kf(x,x)\propto\lVert x\rVert^2$, the variance saturates to the constant
$\satf$. This fixes the divergence \emph{without} normalizing the kernel: $\Hmarg$
cannot grow indefinitely with norm, so to raise utility the optimizer must lower
$\Hcond$ by finding $x^*$ highly correlated with $x_t$; ``super-stimuli'' of
infinite norm yield diminishing returns, embedding the physical constraint that
contrast beyond a point yields no more signal variance.

% ===========================================================================
\subsection{Numerical analysis of the standard utility}
\label{subsec:numerics}

This subsection is entirely about the standard/production utility \eqref{eq:Ustd},
$\Ustd=\Hmarg-\E_{\lgr}[\Hnoise]=I(R;\lgr\mid x^*,\mathcal{D})$, computed by
\texttt{nd\_\allowbreak utility\_\allowbreak new}. (Formulas verified against the live engine; the
\texttt{utility\_\allowbreak numerical\_\allowbreak analysis.tex} line citations from the 2026-05-22 audit
have since drifted, but every formula matches the code.) The spike-count index $r$
below is the model's spike count (Part~I's $\spk$); the count random variable is
$R\sim\Poi(e^{\lgr})$.

\subsubsection{The closed-form decomposition}
With $\lat(x^*)\mid\mathcal{D}\sim\Gauss(\pmean,\pvar)$, $\lgr=\gain\lat+\latz$, so
$\lgr\mid\mathcal{D}\sim\Gauss(\mug,\varg)$, and $R\mid\lgr\sim\Poi(e^{\lgr})$:
\[
  \Hmarg=-\sum_{r=0}^{\infty}p(r)\log p(r),\qquad p(r)=\int\Poi(r;e^{\lgr})\,\Gauss(\lgr;\mug,\varg)\,d\lgr,
\]
the marginal (Poisson--log-normal) entropy of \eqref{eq:Hmarg} (no closed form;
Laplace-approximated below). The subtracted expected-noise entropy \emph{is}
closed-form via the log-normal MGF ($\E[e^{\lgr}]=e^{\mug+\varg/2}$,
$\E[\lgr e^{\lgr}]=e^{\mug+\varg/2}(\mug+\varg)$):
\begin{equation}\label{eq:Ehnoise}
  \E_{\lgr}[\Hnoise]=\underbrace{-e^{\mug+\varg/2}(\mug+\varg-1)}_{\text{exact (MGF), closed form}}\ +\ \underbrace{\sum_{r=0}^{\infty}p(r)\log r!}_{\text{needs }p(r)\text{, Laplace}} ,
\end{equation}
so $\Ustd\approx-\sum_{r=0}^{\rmax}p(r)\log p(r)-\bigl(-e^{\mug+\varg/2}(\mug+\varg-1)+\sum_{r=0}^{\rmax}p(r)\log r!\bigr)$;
both sums truncate at the same $\rmax$. Ground truth: \texttt{nd\_\allowbreak utility\_\allowbreak new}
computes
\begin{lstlisting}
E_H_noise = -torch.exp(mu_g + 0.5*sigma2_g)*(mu_g + sigma2_g - 1) + p_times_logr_sum
\end{lstlisting}

\subsubsection{Laplace approximation and the log-space Lambert-W fix}
The saddle point of $p(r)$ is at the mode $\gmode$ of
$\Poi(r;e^{\lgr})\,\Gauss(\lgr;\mug,\varg)$, solving $r=e^{\gmode}+(\gmode-\mug)/\varg$:
\begin{equation}\label{eq:lambertmode}
  \gmode(r)=r\varg+\mug-\LW\!\bigl(\varg\,e^{\,r\varg+\mug}\bigr),
\end{equation}
$\LW$ the principal Lambert branch. \textbf{The historic bug} (legacy
\texttt{argmax\_\allowbreak g\_\allowbreak old}) formed $\varg e^{\,r\varg+\mug}$ literally and, in float32,
\emph{overflowed} whenever $r\varg+\mug\gtrsim88$ ($e^{88}\sim10^{38}$); masking the
overflowing $r$ let the truncated $\sum p(r)$ exceed $1$ by large factors. This is
the historic ``$\exp(\mu)$ diverges.'' \textbf{The fix} (current
\texttt{\_\allowbreak LambertWLogFunction} $+$ \texttt{\_\allowbreak diff\_\allowbreak argmax\_\allowbreak g}) never forms $e^y$:
set $y=\log\varg+r\varg+\mug$ (finite for finite inputs) and solve
\[
  w+\log w=y\quad\Longleftrightarrow\quad w=\LW(e^y)
\]
by a fixed $10$ Newton iterations directly on the log-form (custom autograd;
backward uses $dW/dy=W/(1+W)$, in which $e^y$ cancels), then
$\gmode(r)=r\varg+\mug-w$. Since $w\approx y-\log y$ for large $y$,
$\gmode(r)\approx-\log\varg+\log y\sim\log r$, so $e^{\gmode(r)}\approx r$ (exactly
right: the Poisson likelihood is maximized when rate $=$ count) and stays in
float32 range for any $r$ --- so $\exp(\gmode)$ is no longer an overflow risk
post-fix.

\subsubsection{Four independent residual failure modes}
Distinct thresholds and remedies (they coincide only because all worsen at large
$\mug/\varg$); (a) and (b) are the pair older docs conflate.

\paragraph{(a) Laplace approximation error in $\log p(r)$.} A local quadratic
model of the integrand; inaccurate when it departs from a quadratic-log shape.
Two regimes: $\varg\to0$ (prior degenerates, saddle singular --- the code falls
back to \emph{exact} Poisson $\log p=\mug r-e^{\mug}-\log r!$ at $\varg<10^{-6}$);
$\varg$ large (flat prior; the asymmetric Poisson factor skews the integrand).
\emph{Independent of $\rmax$} --- each $p(r)$ has its own bias, observable as
$\sum p(r)\neq1$ even with no missing mass (up to $1.015$; that is Laplace error,
not truncation). Measured (vs \texttt{scipy.quad}): error $<1\%$ for $\varg\lesssim1$
at every $\mug$, growing to $5$--$20\%$ at the peak by $\varg\sim50$ (worst where
the predictive peak slides onto $r=0$). The driver is $\varg$, not $\mug$; float32
vs float64 shifts $\log p(r)$ by $\le3.7\times10^{-4}$ nats. \emph{Physical gate:}
an RGC cannot fire more than $\sim100$ spikes in the window, so the sensible regime
is $\varg\lesssim6$, within which (a) is a few-percent effect.

\paragraph{(b) Truncation at $\rmax$.} If the predictive puts non-negligible mass
above $\rmax$, both $-p(r)\log p(r)$ and $p(r)\log r!$ are clipped, biasing
\emph{both} $\Hmarg$ and the closed-sum term low. The moments are
\[
  \E[R]=e^{\mug+\varg/2},\qquad
  \mathrm{Var}(R)=e^{\mug+\varg/2}+e^{2\mug+\varg}(e^{\varg}-1),
\]
the second (log-normal-of-rate) term dominating once $\varg\gtrsim1$. As $\varg$
grows the predictive is strongly right-skewed and its centres separate: the mean
$e^{\mug+\varg/2}$ climbs, the median stays $\approx e^{\mug}$, the mode drifts
\emph{down} toward $0$ ($\approx e^{\mug-\varg}$) --- so ``$p(r)$ peaks at
$e^{\mug}$'' is only the small-$\varg$ limit. A sufficient benign-truncation
condition is $\rmax\gtrsim\E[R]+k\sqrt{\mathrm{Var}(R)}$. The code
\texttt{compute\_\allowbreak adaptive\_\allowbreak rmax} is stricter (a $k$-$\sigma$ upper tail in
$\lgr$-space, exponentiated to a rate, plus a Poisson std, plus a margin):
\[
  \rmax^{\mathrm{adapt}}=e^{\mug+k\sigma_g}+5\sqrt{e^{\mug+k\sigma_g}}+10,\quad k=3,
\]
clamped to $[200,10000]$ with a \texttt{ValueError} if the needed $\rmax$ exceeds
the ceiling (never silently over-truncates; a further guard raises if
$\text{max\_\allowbreak rmax}>e^{80}$ so the overflow clamp cannot hide failures). The fixed
default $\rmax=100$ is benign for confident predictions but \emph{not safe in
general}: it fails whenever $\mug\gtrsim\log100\approx4.6$ \emph{or} $\varg$ grows
large --- exactly the high-uncertainty regime active learning explores.

\paragraph{(c) Float32 overflow of the exact closed-form term.} The exact term
$-e^{\mug+\varg/2}(\mug+\varg-1)$ overflows float32 when
$\mug+\tfrac12\varg\gtrsim88$, sending $\E_{\lgr}[\Hnoise]\to-\infty$ and
$\Ustd\to+\infty$. This is the \emph{only} residual ``$\exp$ can diverge'' issue,
and it lives in an \emph{exact analytical} term, not the Laplace approximation ---
distinct from the historic bug. In practice truncation (b) bites long first
($\sigma_g\sim1\Rightarrow$ truncation near $\mug\sim1$ at $\rmax=100$, vs
$\mug\sim88$ here).

\paragraph{(d) Catastrophic cancellation.} Inherent to the decomposition (survives
even with exact Laplace and no truncation). \emph{Inside} $\E_{\lgr}[\Hnoise]$
(large, earliest): both $-e^{\mug+\varg/2}(\mug+\varg-1)$ and $\sum p(r)\log r!$
grow $\sim e^{\mug}$, while the true $\E_{\lgr}[\Hnoise]$ grows only \emph{linearly}
($\sim\tfrac12\mug$); at $\mug=10$ each term $\sim2\times10^5$ vs a true value
$\sim6$ ($\sim4.5$ digits cancel), at $\mug=15$ $\sim6.7$ digits cancel ---
essentially noise in float32's $\sim7$-digit mantissa (and truncating the second
sum makes this \emph{worse}). \emph{Between} $\Hmarg$ and $\E_{\lgr}[\Hnoise]$
(smaller, second-level): by Jensen $\Hmarg\ge\E_{\lgr}[\Hnoise]$ while both grow
linearly in $\mug$, a second cancellation costing fewer digits. Unfixable without
rewriting the decomposition so cancelling pieces are computed together (e.g.\
summing $-p(r)\log p(r)-p(r)\log r!$ as a single term before any large analytical
exponential). Flagged, out of scope.

\subsubsection{``Utility diverges under gradient ascent,'' re-examined}
The claim is \emph{two} distinct things. (i) \textbf{Kernel-driven input blow-up:}
$\Kf(\scal x,\scal x)=\scal^2\Kf(x,x)\Rightarrow$ ascending $\Ustd$ with no norm
constraint sends $\pvar\propto\lVert x\rVert^2\to\infty$, hence $\varg\to\infty$
(the mechanism of \S\ref{subsec:divergence}--\ref{subsec:kernelfixes}). (ii) The
\textbf{computed} utility then misbehaves via the four numerical modes, typically
in order: truncation collapse of $\Hmarg$ (b), then Laplace error at large $\varg$
(a), then cancellation inside $\E_{\lgr}[\Hnoise]$ (d) as $\mug$ grows, finally
float32 overflow (c) at extreme scales. The kernel makes the inputs unbounded;
then numerics make the computed utility ill-behaved --- not the same thing, and
the four numerical effects are not the same thing either. Numerical fixes alone
will not cure it: either constrain the optimizer (norm projection, or the
normalized kernel of Proposition~\ref{prop:normkernel}) or rewrite the closed-form
term in an asymptotically stable form.

\caveat{Utility-accuracy ceiling (known limitation). The standard utility can run
to $+\infty$ and the DA utility can silently collapse, both via the $\rmax=100$
truncation, at high firing; the production loop's $f_{\max}$ guard keeps generation
in the accurate regime, and a DA-path lookup-table was deliberately not built.
This is a documented limitation, not a resolved result.}
